% SGA 6, Exposición IV: fragmento español alineado con la fuente.
% Parte 2 de 7. Alcance: sección 2.1 hasta (2.6.2), líneas 6209--6303 inclusive.
% Fuente francesa: sga6_fr_r1/src/fr.tex.
% SHA-256 de la fuente francesa completa:
% 0D2158B9FE3B466309A92165727A0B191422E05FBAA0DACD5806E8322C54BA1A.
\providecommand{\sisi}[2]{#2}

\section*{2. Los funtores \(K_{\bullet}\) y \(K'\) de un topos anillado}

\textbf{2.1.} Sea \((X,A)\) un topos anillado. Recordemos las notaciones siguientes:

\(D^b({}_A X)_{\coh}\) = categoría de los complejos pseudocoherentes y de cohomología acotada de \(A\)-módulos a izquierda (I 2.15).

\(D({}_A X)_{\parf}\) = categoría de los complejos de \(A\)-módulos a izquierda de amplitud perfecta finita (I 5.8.4).

\(\Loclib({}_A X)\) = categoría de los \(A\)-módulos a izquierda que son localmente sumandos directos de módulos libres de tipo finito (II 2.0).

\(\Coh({}_A X)\) = categoría de los \(A\)-módulos a izquierda coherentes (con respecto a la categoría de los \(A\)-módulos libres de tipo finito) (I 3.1).

Las dos últimas categorías son subcategorías plenas de la categoría \(\Mod({}_A X)\) de los \(A\)-módulos a izquierda, las dos primeras subcategorías trianguladas plenas de \(D({}_A X)\) (\(= D(\Mod({}_A X))\)).

\textbf{Definición 2.2.} Sea \((X,A)\) un topos anillado. Salvo mención en contrario, pondremos
\[
 \begin{gathered}
  K_\bullet({}_A X)=k(D^b({}_A X)_{\coh}),\\
  K'({}_A X)=k(D({}_A X)_{\parf}),
 \end{gathered}
\]
donde \(k(\ )\) está definido en (1.1), y
\[
 \begin{gathered}
  K_\bullet({}_A X)_{\naif}=k(\Coh({}_A X)),\\
  K'({}_A X)_{\naif}=k(\Loclib({}_A X)),
 \end{gathered}
\]
donde \(k(\ )\) está definido por (1.5.2), siendo la categoría abeliana ambiente \(\Mod({}_A X)\)\footnote{Nos permitiremos omitir el índice \(A\) cuando de ello no pueda resultar confusión. A veces también suprimiremos el índice «~ingenuo~», pero indicándolo explícitamente.}. Diremos que \(K_\bullet({}_A X)\) (resp. \(K'({}_A X)\)) es el grupo de Grothendieck de los complejos pseudocoherentes (resp. perfectos) de \(A\)-módulos a izquierda, y que \(K_\bullet({}_A X)_{\naif}\) (resp. \(K'({}_A X)_{\naif}\)) es el grupo de Grothendieck de los \(A\)-módulos a izquierda coherentes (resp. localmente libres).

\textbf{Observaciones 2.3.}
\begin{enumerate}[label=\alph*)]
\item Se pueden definir variantes de \(K_\bullet(\ )\) y \(K'(\ )\) poniendo
\[
 K_\bullet^{\loc}({}_A X)=k(D^{\loc b}({}_A X)_{\coh})
\]
(donde «~\(\loc b\)~» significa «~de cohomología localmente acotada~»), y
\[
 K_{\loc}'({}_A X)=k(D({}_A X)_{\parf}).
\]
Estos grupos coinciden con los definidos en (2.2) cuando el objeto final de \(X\) es cuasicompacto, pero en general se prestan peor al cálculo de las intersecciones.
\item El grupo \(K_\bullet({}_A X)_{\naif}\) prácticamente solo tiene interés si \(A\) es coherente (de otro modo existe el riesgo de que haya muy pocos haces coherentes \ldots).
\item La consideración del \(k\) de categorías trianguladas «~demasiado grandes~» no ofrece interés. Así, de (1.2 b)) se deducen fácilmente las relaciones
\[
 k(D({}_A X))=k(D^-({}_A X))=k(D^+({}_A X))=k(D^b({}_A X))=k(D^-({}_A X)_{\coh})=0
\]
(para la última igualdad, obsérvese que \(L\in\Ob D^-({}_A X)_{\coh}\) implica \(\coprod_{i\geq0}L[2i]\in\Ob D^-({}_A X)_{\coh}\)).
\end{enumerate}

\textbf{2.4. Caso de un haz de anillos coherente.} Sea \((X,A)\) un topos anillado, donde se supone que \(A\) es coherente como \(A\)-módulo a izquierda (I 3.1 y 3.6). Todo \(A\)-módulo a izquierda coherente es un complejo pseudocoherente (I 3.5), por lo que se tiene un funtor plenamente fiel
\begin{equation}\tag{2.4.1}
 \Coh({}_A X)\hookrightarrow D^b({}_A X)_{\coh}.
\end{equation}
Ahora bien, (I 3.4) \(\Coh({}_A X)\) es una subcategoría gruesa de \(\Mod({}_A X)\), y, en virtud de (I 3.5), \(D^b({}_A X)_{\coh}\) es la subcategoría plena de \(D^b({}_A X)\) formada por los objetos \(L\) tales que \(H^i(L)\in\Ob\Coh({}_A X)\) para todo \(i\). Por consiguiente, se deduce de (1.6) que el homomorfismo
\begin{equation}\tag{2.4.2}
 K_\bullet({}_A X)_{\naif}\longrightarrow K_\bullet({}_A X)
\end{equation}
definido por (2.4.1) es un isomorfismo.

\textbf{2.5. El homomorfismo \(K'\to K_\bullet\).} Sea \((X,A)\) un topos anillado. El funtor de inclusión \(D({}_A X)_{\parf}\hookrightarrow D^b({}_A X)_{\coh}\) define un homomorfismo
\begin{equation}\tag{2.5.1}
 K'({}_A X)\longrightarrow K_\bullet({}_A X).
\end{equation}
Este no es en general ni inyectivo ni sobreyectivo. No obstante, hagamos las hipótesis siguientes:
\begin{enumerate}[label=(\roman*)]
\item el objeto final de \(X\) es cuasicompacto;
\item \(X\) tiene suficientes puntos, y, para todo punto \(x\) de \(X\), todo \(A_x\)-módulo a izquierda es de dimensión tor finita.
\end{enumerate}
Entonces se tiene, según (I 5.10), \(D({}_A X)_{\parf}=D^b({}_A X)_{\coh}\), y el homomorfismo (2.5.1) es un isomorfismo. Las hipótesis (i) y (ii) se satisfacen, por ejemplo, en el caso de un espacio compacto anillado por anillos locales regulares.

\textbf{2.6. Emparejamientos entre \(K'\) y \(K'\) (caso no conmutativo).} Sea
\[
 f:(X',A')\longrightarrow (X,A)
\]
un morfismo de topoi anillados. Denotemos por \(\Mod({}_{A'} X'_A)\) la categoría de los bimódulos, a izquierda sobre \(A'\), a derecha sobre \(f^{-1}(A)\), y por \(\Loclib_{A'}({}_{A'} X'_A)\) la subcategoría plena formada por los objetos que pertenecen a \(\Loclib({}_{A'} X')\). Por analogía con (2.2), pondremos
\[
 K'_{\naif}({}_{A'} X'_A)_{A'}=k(\Loclib_{A'}({}_{A'} X'_A)).
\]
El producto tensorial
\[
 \otimes_A:\Mod({}_{A'} X'_A)\times\Mod({}_A X)\longrightarrow \Mod({}_{A'} X')
\]
envía
\[
 (E',E)\longmapsto E'\otimes_{f^{-1}(A)}f^{-1}(E),
\]
e induce un funtor
\begin{equation}\tag{2.6.1}
 \Loclib_{A'}({}_{A'} X'_A)\times\Loclib({}_A X)\longrightarrow\Loclib({}_{A'} X').
\end{equation}
La aplicación \((E',E)\mapsto\cl(E'\otimes_A E)\) deducida de (2.6.1) es biaditiva en el sentido de (1.5), y define por tanto un homomorfismo
\begin{equation}\tag{2.6.2}
 K'_{\naif}({}_{A'} X'_A)_{A'}\otimes K'_{\naif}({}_A X)\longrightarrow K'_{\naif}({}_{A'} X').
\end{equation}

